\documentclass[12pt]{article}

% Codificación e idioma
\usepackage[utf8]{inputenc}
\usepackage[T1]{fontenc}
\usepackage[spanish,mexico]{babel}

% Matemáticas
\usepackage{amsmath,amssymb}

% Formato y utilidades
\usepackage{geometry}
\usepackage{enumitem}
\usepackage{booktabs}
\usepackage{longtable}
\usepackage{array}
\usepackage{xcolor}
\usepackage{hyperref}
\usepackage{graphicx}
\usepackage{fancyhdr}

% Márgenes
\geometry{letterpaper,margin=2.5cm}

% Listas
\setlist[itemize]{leftmargin=1.5em}
\setlist[enumerate]{leftmargin=1.8em}

% Enlaces
\hypersetup{
    colorlinks=true,
    linkcolor=blue,
    urlcolor=blue,
    citecolor=blue
}

% Encabezados y pies de página
\pagestyle{fancy}
\fancyhf{}

% ---------- Encabezado ----------
\fancyhead[L]{Tecnológico Nacional de México}
\fancyhead[R]{Cálculo Integral}

% ---------- Pie de página ----------
\fancyfoot[L]{Carlos Ernesto Martínez Rodríguez}
\fancyfoot[C]{\thepage}
\fancyfoot[R]{mrdgzcarlose@gmail.com}

\renewcommand{\headrulewidth}{0pt}
\renewcommand{\footrulewidth}{0.4pt}

% Título y subtítulo
\title{
    \textbf{Cálculo Integral}\\[0.5em]
    \large\textit{Tecnológico Nacional de México}
}

\author{Carlos E. Martínez Rodríguez}
\date{}

\begin{document}

\maketitle
\tableofcontents

\section*{Información del Curso}

\begin{itemize}
    \item \textbf{Asignatura:} Cálculo Integral.
    \item \textbf{Clave:} ACF--0902.
    \item \textbf{Horas teoría--práctica--créditos:} 3--2--5.
    \item \textbf{Periodo:} Agosto--Diciembre 2026.
\end{itemize}

% ============================================================
% CÓDIGO QR DEL TEMARIO
% ============================================================
%
% Coloca el archivo PNG del código QR en la misma carpeta que
% este archivo .tex y sustituye, si es necesario, el nombre
% QR_Temario_Calculo_Integral_TNM.png por el nombre real.
%
% Ejemplo:
%
% \begin{center}
%     \includegraphics[width=0.22\textwidth]
%     {QR_Temario_Calculo_Integral_TNM.png}
%
%     \small
%     Escanea el código QR para consultar el temario del curso.
% \end{center}
%
% ============================================================

\section*{Objetivo general del curso}

El curso tiene como propósito que el estudiante \textbf{aplique la definición de integral y las técnicas de integración para resolver problemas de ingeniería}.

\section{Teorema fundamental del cálculo}

\textbf{Objetivos:} Comprender los dos teoremas fundamentales del cálculo para establecer la relación entre el cálculo diferencial y el cálculo integral, y aplicar los teoremas y propiedades de la integral para evaluar integrales definidas.

\subsection{Subtemas}

\begin{enumerate}
    \item Introducción histórica al cálculo integral.
    \item El problema del área.
    \item Aproximación del área de regiones planas.
    \item Particiones de un intervalo.
    \item Sumas de Riemann.
    \item Aproximaciones por extremo izquierdo.
    \item Aproximaciones por extremo derecho.
    \item Aproximaciones mediante el punto medio.
    \item Definición de integral definida.
    \item Condiciones de existencia de la integral definida.
    \item Propiedades de la integral definida.
    \item Teorema del valor intermedio.
    \item Primer Teorema Fundamental del Cálculo.
    \item Segundo Teorema Fundamental del Cálculo.
    \item Evaluación de integrales definidas.
    \item Interpretación geométrica de la integral definida.
\end{enumerate}

\textbf{Tiempo previsto:} La planeación considera \textbf{20 horas} para esta unidad, distribuidas en 12 horas teóricas y 8 horas prácticas, equivalentes a aproximadamente cuatro semanas.

\section{Integral indefinida y métodos de integración}

\textbf{Objetivos:} Identificar el método de integración más adecuado para resolver una integral indefinida y utilizar diferentes técnicas de integración en funciones algebraicas, trigonométricas, exponenciales y logarítmicas.

\subsection{Subtemas}

\begin{enumerate}
    \item Concepto de función primitiva o antiderivada.
    \item Integral indefinida.
    \item Constante de integración.
    \item Propiedades de la integral indefinida.
    \item Integración directa.
    \item Cambio de variable o sustitución.
    \item Integración por partes.
    \item Integrales de funciones trigonométricas.
    \item Sustituciones trigonométricas.
    \item Integración de funciones racionales.
    \item Descomposición en fracciones parciales.
    \item Integrales de funciones exponenciales.
    \item Integrales de funciones logarítmicas.
    \item Selección del método de integración adecuado.
    \item Uso de herramientas computacionales para verificar integrales.
\end{enumerate}

\textbf{Tiempo previsto:} La planeación considera \textbf{20 horas} para esta unidad, distribuidas en 12 horas teóricas y 8 horas prácticas, equivalentes a aproximadamente cuatro semanas.

\section{Aplicaciones de la integral}

\textbf{Objetivos:} Utilizar las definiciones de integral y las técnicas de integración para resolver problemas geométricos y problemas aplicados en ingeniería.

\subsection{Subtemas}

\begin{enumerate}
    \item Área bajo una curva.
    \item Área entre curvas.
    \item Volúmenes de sólidos de revolución.
    \item Método de discos y arandelas.
    \item Método de cascarones cilíndricos.
    \item Longitud de arco.
    \item Área de superficies de revolución.
    \item Trabajo.
    \item Centro de masa.
    \item Integrales impropias.
    \item Aplicaciones de la integral en problemas de ingeniería.
    \item Uso de software para plantear, resolver y verificar problemas de aplicación.
\end{enumerate}

\textbf{Tiempo previsto:} La planeación considera \textbf{20 horas} para esta unidad, distribuidas en 12 horas teóricas y 8 horas prácticas, equivalentes a aproximadamente cuatro semanas.

\section{Series}

\textbf{Objetivos:} Aplicar sucesiones y series para analizar procesos de convergencia y utilizar series de potencias, de Taylor y de Maclaurin para aproximar funciones e integrales especiales.

\subsection{Subtemas}

\begin{enumerate}
    \item Sucesiones.
    \item Series finitas e infinitas.
    \item Convergencia y divergencia.
    \item Criterios de convergencia.
    \item Criterios de divergencia.
    \item Series de potencias.
    \item Intervalo de convergencia.
    \item Radio de convergencia.
    \item Derivación de series de potencias.
    \item Integración de series de potencias.
    \item Serie de Taylor.
    \item Serie de Maclaurin.
    \item Representación de funciones mediante series.
    \item Aproximación de funciones mediante polinomios de Taylor.
    \item Evaluación aproximada de integrales mediante series de Taylor.
    \item Uso de software para visualizar aproximaciones mediante series.
\end{enumerate}

\textbf{Tiempo previsto:} La planeación considera \textbf{20 horas} para esta unidad, distribuidas en 12 horas teóricas y 8 horas prácticas, equivalentes a aproximadamente cuatro semanas.

\section{Planeación}

\begin{center}
\begin{tabular}{>{\raggedright\arraybackslash}p{8cm}c}
\toprule
\textbf{Unidad} & \textbf{Horas previstas} \\
\midrule
1. Teorema fundamental del cálculo & 20 \\
2. Integral indefinida y métodos de integración & 20 \\
3. Aplicaciones de la integral & 20 \\
4. Series & 20 \\
\midrule
\textbf{Total} & \textbf{80} \\
\bottomrule
\end{tabular}
\end{center}

\section{Criterios generales de evaluación}

En cada unidad se establece, de manera general, la siguiente ponderación:

\begin{itemize}
    \item \textbf{70\%}: evaluación escrita y presencial.
    \item \textbf{30\%}: actividades de aprendizaje a criterio del docente.
\end{itemize}

Las actividades de enseñanza y aprendizaje contemplan, entre otras, la participación activa en clase, trabajo individual y en equipo, resolución de ejercicios, trabajo extra-clase, uso de herramientas computacionales, atención de dudas y evaluación escrita.

\section{Bibliografía básica}

\begin{itemize}
    \item Stewart, James. \textit{Cálculo de una variable}. 8a. edición. Cengage Learning.
    \item Larson, Ron y Hostetler, Robert. \textit{Cálculo}. 8a. edición. McGraw-Hill.
    \item Leithold, Louis. \textit{El Cálculo}. Oxford University Press.
    \item Purcell, Edwin J. y Dale Varberg. \textit{Cálculo}. Prentice Hall.
    \item Zill, Dennis G. \textit{Cálculo: trascendentes tempranas}. 4a. edición. McGraw-Hill.
\end{itemize}

\newpage

\section{Material complementario del curso}

% ============================================================
% AMBIENTE PARA INSERTAR CÓDIGOS QR DE NOTAS O MATERIAL
% ============================================================
%
% Ejemplo de uso:
%
% \begin{center}
%     \includegraphics[width=0.22\textwidth]{QR_Notas_CI_Unidad1.png}
%
%     \small
%     Escanea el código QR para acceder a las notas de la Unidad 1.
% \end{center}
%
% Puedes repetir este bloque para cada unidad o documento.
%
% ============================================================

\end{document}
